\newpage{\pagestyle{empty}\cleardoublepage}

\newpage
\begin{center}
\thispagestyle{empty} \vspace*{0.5cm} \textbf{\huge
Emisión de paquetes de n-fonones en una molécula excitónica mediante procesos de Stokes}\\[3.0cm]
\Large\textbf{Jhon Sebastian García Barrera}\\[3.5cm]
\small Trabajo de grado presentado como requisito para optar al
t\'{\i}tulo de:\\
\textbf{Físico}\\[3.0cm]
Director:\\
Ph.D. Edgar A. Gómez G.\\
Codirector:\\
Ph.D. Santiago Echeverri Arteaga\\[3.0cm]
Universidad del Quindío\\
Facultad de Ciencias Básicas y Tecnologías\\
Programa de Física\\
Armenia, Colombia\\
2026\\
\end{center}

\newpage{\pagestyle{empty}\cleardoublepage}

\newpage
\thispagestyle{empty} \textbf{}\normalsize
\\\\\\%
[4.0cm]

\begin{flushright}
\begin{minipage}{10cm}
    \noindent
        \small
        La sabiduría no se puede transmitir. La sabiduría que un sabio intenta comunicar siempre suena a locura para otro. El conocimiento se puede comunicar, pero la sabiduría no. Se puede encontrar, vivir, experimentar. Se puede expresar, pero no enseñar. Por eso cada camino es propio, y cada búsqueda verdadera exige tiempo, paciencia y fidelidad a uno mismo.\\\\
        \textbf{Hermann Hesse}\\
\end{minipage}
\end{flushright}

\newpage{\pagestyle{empty}\cleardoublepage}

\newpage
\thispagestyle{empty} \textbf{}\normalsize
\\\\\\%
\textbf{\LARGE Agradecimientos}
\addcontentsline{toc}{chapter}{\numberline{}Agradecimientos}\\\\

A mis directores, profesores, compañeros y amigos, les debo algo que no aparece en ninguna ecuación de este trabajo: la forma en que aprendí a pensar. Cada conversación, cada corrección, cada pregunta sin respuesta fácil dejó una huella que no siempre supe reconocer en el momento, pero que está presente en cada página de este texto.

A mis padres y familiares, les debo algo de un tipo diferente y más difícil de expresar. Probablemente son los últimos en reconocer cuánto han contribuido a este trabajo, no solo con su apoyo, sino permitiéndome seguir adelante incluso cuando el camino fue largo. Quien haya emprendido un proyecto como este sabe lo que eso significa y lo que, a veces, cuesta. No sé cómo darles las gracias.

\newpage{\pagestyle{empty}\cleardoublepage}

\newpage
\addcontentsline{toc}{chapter}{\numberline{}Resumen}
\textbf{\LARGE Resumen}\\

Se estudia la generación de paquetes de $n$ fonones en una molécula excitónica formada por dos puntos cuánticos vinculados mediante interacción de Förster y acoplados a una nanocavidad acústica de THz. El acoplamiento de Förster genera una base de estados colectivos brillante y oscuro que introduce un factor superradiante $\sqrt{2}$ en la frecuencia super-Rabi efectiva y un corrimiento $-J$ en las resonancias de Stokes. Se derivan analíticamente las condiciones de resonancia y las frecuencias efectivas en tres regímenes de parámetros, y se resuelve numéricamente la ecuación maestra de Lindblad con QuTiP para calcular funciones de correlación de orden superior, trayectorias cuánticas de Monte Carlo y mapas de pureza $\Pi_n$ para $n=2,3,4$. Se muestra que la estadística de los paquetes emitidos es sintonizable entre los regímenes de n-phonon gun, n-phonon laser y estado térmico de paquetes de paquetes ajustando la relación $\gamma/\kappa$. Los resultados muestran que para $\lambda/\omega_b \approx 0.22$ se obtienen $\Pi_2 \approx 98.6\%$ y $\Pi_3 \approx 98.9\%$, y que $\Pi_3 > 95\%$ es accesible con factores de calidad $Q \sim 90$--$300$, frente a $Q \sim 500$ requeridos en la literatura, constituyendo una propuesta experimentalmente viable con la tecnología actual.

\vspace{0.5cm}
\textbf{\small Palabras clave: } fonónica cuántica, paquetes de fonones, molécula excitónica, interacción de Förster, resonancia de Stokes, pureza de emisión, superradiancia. \\

\textbf{\LARGE Abstract}\\

The generation of $n$-phonon bundles in an excitonic molecule formed by two quantum dots coupled via Förster interaction and coupled to a THz acoustic nanocavity is studied. The Förster coupling generates a basis of collective bright and dark states that introduces a superradiant factor $\sqrt{2}$ in the effective super-Rabi
frequency and a shift $-J$ in the Stokes resonances. Resonance conditions and effective frequencies are derived analytically in three parameter regimes, and the Lindblad master equation is solved numerically with QuTiP to compute higher-order correlation functions, quantum Monte Carlo trajectories, and emission purity maps $\Pi_n$ for $n=2,3,4$. The bundle statistics are shown to be tunable between single-bundle source, laser, and thermal bundle regimes by adjusting the ratio $\gamma/\kappa$. The results show that for $\lambda/\omega_b \approx 0.22$, $\Pi_2 \approx 98.6\%$ and $\Pi_3 \approx 98.9\%$ are achieved, and that $\Pi_3 > 95\%$ is accessible with quality factors $Q \sim 90$--$300$, compared to $Q \sim 500$ required in the literature, constituting an experimentally viable proposal with current technology.

\vspace{0.5cm}
\textbf{\small Key words: } quantum phononics, phonon bundles, excitonic molecule, Förster interaction, Stokes resonance, emission purity, superradiance.